% Part: history
% Chapter: biographies 
% Section: alfred-tarski

\documentclass[../../../include/open-logic-section]{subfiles}

\begin{document}

\olfileid{his}{bio}{tar}

\olsection{ఆల్ఫ్రెడ్ టార్స్కీ}

\olphoto{tarski-alfred}{ఆల్ఫ్రెడ్ టార్స్కీ}

ఆల్ఫ్రెడ్ టార్స్కీ 1901 జనవరి 14న పోలండ్‌లోని వార్సాలో (అప్పుడు రష్యన్ సామ్రాజ్యంలో భాగం) జన్మించారు. ``నెపోలియన్‌లాంటి'' వ్యక్తిగా వర్ణించబడిన టార్స్కీ ఉత్సాహంగా, ఎక్కువగా మాట్లాడే, తీవ్ర స్వభావం కలవారు. ఆ శక్తి ఆయన ఉపన్యాసాలలో తరచూ కనిపించేది---ఒకసారి ఉపన్యాసం మధ్యలో సిగరెట్ పారేస్తూ చెత్తబుట్టకు నిప్పంటించారు; ఆ తరువాత ఆ భవనంలో ఉపన్యాసాలివ్వకూడదని నిషేధించారు.

చిన్నతనం నుంచే టార్స్కీకి జ్ఞానదాహం ఉండేది. తరువాతి కాలంలో, తర్కశాస్త్రంలో మాత్రమే బీ గ్రేడ్ వచ్చిందని అందుకే దాన్ని చదివానని విద్యార్థులకు చెప్పేవారు; కానీ ఉన్నత పాఠశాల రికార్డులు తర్కశాస్త్రంలో కూడా, అన్ని విషయాల్లోనూ, ఏ గ్రేడే వచ్చిందని చూపుతున్నాయి. 1918 నుంచి 1924 వరకు వార్సా విశ్వవిద్యాలయంలో చదివారు. మొదట జీవశాస్త్రం చదవాలనుకున్నా, ఆ విశ్వవిద్యాలయం వార్సా తర్కశాస్త్ర-తత్వశాస్త్ర పాఠశాలకు కేంద్రం కావడంతో గణితం, తత్వశాస్త్రం, తర్కశాస్త్రంపై ఆసక్తి పెంచుకున్నారు. స్టానిస్లావ్ లెష్నియెవ్స్కీ పర్యవేక్షణలో 1924లో డాక్టరేట్ పొందారు.

1939లో అమెరికా సంయుక్త రాష్ట్రాలకు వలస వెళ్లే ముందు, వార్సాలో మాధ్యమిక పాఠశాల ఉపాధ్యాయునిగా పనిచేస్తూనే టార్స్కీ తన అతి ముఖ్యమైన కృషిలో కొంత పూర్తి చేశారు. తార్కిక అనుగమనం, తార్కిక సత్యంపై ఆయన రచనలు ఆ కాలంలోనే వచ్చాయి. 1939లో ఉపన్యాసాల పర్యటన కోసం టార్స్కీ అమెరికా సంయుక్త రాష్ట్రాల్లో ఉన్నారు. ఆయన అక్కడ ఉన్నప్పుడే జర్మనీ పోలండ్‌పై దండెత్తింది; యూదు మూలాలు ఉండడంతో టార్స్కీ తిరిగి వెళ్లలేకపోయారు. యుద్ధం ముగిసే వరకు ఆయన భార్య, పిల్లలు పోలండ్‌లోనే ఉన్నారు; తరువాత వారు కూడా అమెరికాకు వలస రాగలిగారు. టార్స్కీ హార్వర్డ్‌లో, న్యూయార్క్ సిటీ కాలేజీలో, ప్రిన్స్‌టన్‌లోని ఇన్‌స్టిట్యూట్ ఫర్ అడ్వాన్స్‌డ్ స్టడీలో, చివరకు బర్కిలీలోని కాలిఫోర్నియా విశ్వవిద్యాలయంలో బోధించారు. అక్కడ తర్కశాస్త్రం, విజ్ఞానశాస్త్ర పద్ధతిశాస్త్రం పై బహుశాఖా కార్యక్రమాన్ని స్థాపించారు. 1983 అక్టోబరు 26న 82 ఏళ్ల వయసులో మరణించారు.

\begin{reading}
టార్స్కీ జీవితం గురించి మరిన్ని వివరాలకు \emph{Alfred Tarski: Life
  and Logic} అనే జీవిత చరిత్ర \citep{Feferman2004} చూడండి. తార్కిక అనుగమనం, సత్యంపై ఆయన ప్రాముఖ్యమైన రచనలు ఆంగ్లంలో \citep{Tarski1983}లో లభిస్తాయి. టార్స్కీ మూల రచనలన్నీ నాలుగు సంపుటాల శ్రేణిగా సంకలితం అయ్యాయి \citep{Tarski1981}.
\end{reading}

\end{document}
